% Part: history
% Chapter: biographies
% Section: ernst-zermelo

\documentclass[../../../include/open-logic-section]{subfiles}

\begin{document}

\olfileid[id]{his}{bio}{zer}

\olsection{Ernst Zermelo}

\olphoto{zermelo-ernst}{Ernst Zermelo}

Ernst Zermelo lahir pada 27 Juli 1871 di Berlin, Jerman. Ia mempunyai lima
saudari, tetapi keluarganya menderita kesehatan yang buruk dan hanya tiga
dari mereka bertahan hingga dewasa. Orang tuanya juga meninggal ketika ia
masih muda, sehingga ia dan saudara-saudaranya menjadi yatim piatu ketika ia
berusia tujuh belas tahun. Zermelo sangat berminat pada seni, terutama puisi. Ia
dikenal tajam, jenaka, dan kritis. Pencapaian matematisnya yang paling
terkenal mencakup pengenalan aksioma pilihan (pada 1904) dan aksiomatisasi
teori himpunannya (pada 1908).

Minat Zermelo di universitas beragam. Ia mengikuti kuliah fisika, matematika,
dan filsafat. Di bawah bimbingan Hermann Schwarz, Zermelo menyelesaikan
disertasinya, \emph{Penyelidikan dalam Kalkulus Variasi}, pada 1894 di
Universitas Berlin. Pada 1897, ia memutuskan melanjutkan studi di Universitas
G\"{o}ttingen, tempat ia sangat dipengaruhi oleh karya landasan David
Hilbert. Pada 1899, ia memenuhi syarat menjadi profesor, tetapi baru
memperoleh jabatan tersebut sebelas tahun kemudian---mungkin karena sikapnya
yang aneh dan ``ketergesaan gugup''.

Zermelo akhirnya memperoleh jabatan profesor bergaji di Universitas Zurich
pada 1910, tetapi terpaksa pensiun pada 1916 karena tuberkulosis. Setelah
pulih, ia diberi jabatan profesor kehormatan di Universitas Freiburg pada
1921. Pada masa ini ia mengerjakan landasan matematika. Ia merasa terganggu
oleh karya Thoralf Skolem dan Kurt G\"{o}del, dan secara terbuka mengkritik
pendekatan mereka dalam makalah-makalahnya. Ia diberhentikan dari jabatannya
di Freiburg pada 1935 karena ketidakpopulerannya dan penentangannya terhadap
kebangkitan Hitler di Jerman.

Tahun-tahun terakhir kehidupan Zermelo ditandai oleh keterasingan. Setelah
diberhentikan pada 1935, ia meninggalkan matematika. Ia pindah ke pedesaan
dan hidup sederhana. Ia menikah pada 1944 dan menjadi sepenuhnya bergantung
pada istrinya ketika penglihatannya memburuk. Zermelo kehilangan seluruh
penglihatannya pada 1951. Ia meninggal di G\"{u}nterstal, Jerman, pada
21 Mei 1953.

\begin{reading}
Untuk biografi lengkap Zermelo, lihat \citet{Ebbinghaus2015}. Makalah
penting Zermelo dari 1904 dan 1908 tersedia dalam bahasa Jerman asli
\citep{Zermelo1904,Zermelo1908}. Kumpulan karya Zermelo, termasuk
tulisannya mengenai fisika, tersedia dalam terjemahan bahasa Inggris dalam
\citep{Ebbinghaus2010,Ebbinghaus2013}.
\end{reading}

\end{document}
